\chapter{Introduction}
\label{ch:introduction}

\section{The problem: hybrid serving is controlled one layer at a time}

Multi-tenant platforms no longer serve one kind of traffic. The same tenants,
on the same clusters, issue classic CRUD requests and AI inference requests
--- chat completions, embeddings, multi-step agent chains --- often within a
single user-facing transaction. Three control surfaces govern what that
serving costs and how it feels:

\begin{enumerate}
  \item \textbf{Replica autoscaling} decides how much capacity each service
        holds. In production this is the Kubernetes Horizontal Pod Autoscaler
        (HPA)~\citep{k8shpa}, KEDA~\citep{keda}, or a research
        controller~\citep{firm2020,sageserve2025,chiron2025,aladdin2024}.
  \item \textbf{Semantic response caching} decides which AI answers are
        served from an embedding-similarity cache instead of a model
        call~\citep{gptcache2023,scalm2024,instcache2024,meancache2024}.
  \item \textbf{Model-tier routing} decides which of several differently
        priced model classes answers each request~\citep{infaas2019}.
\end{enumerate}

Each layer has a mature literature and, in every surveyed system, an
independent controller: an autoscaler that has never heard of the cache, a
cache that does not know what a replica costs, a router that optimizes one
request at a time. Each controller can be locally optimal. This thesis starts
from the observation that the three layers share one budget, one latency
target, and one set of tenants --- and asks whether local optima compose.

\section{Thesis statement}

\begin{quote}
\emph{Per-layer local optima do not compose into a joint optimum. A
controller that plans replicas, semantic-cache budget, and model tier
together --- one receding-horizon optimization per control interval, against
a single objective in cost, SLO compliance, and cross-tenant fairness ---
dominates tuned per-layer baselines on that objective, and does so with
per-tenant dollar accounting and budget guardrails that no surveyed system
provides.}
\end{quote}

The claim is tested, not asserted: against five baselines each grid-tuned on
the thesis's own objective, under pre-registered protocols (seven at the
first campaign, forty-eight by the last) whose hypotheses --- including the
ones that failed --- were committed and pushed before the first run.

\section{What was built}

\polyforge{} is a Kubernetes-native control plane implementing the thesis
statement end to end (Chapter~\ref{ch:design}):

\begin{itemize}
  \item a \textbf{control plane} (Go) for tenants, projects, API keys, rate
        limits, and telemetry, with row-level-security-enforced tenant
        isolation proven by test;
  \item an \textbf{AI gateway} (Go) with multi-backend chat routing, a
        semantic cache with cost-aware eviction, vector search, and an agent
        loop;
  \item a \textbf{workload classifier} mapping 12-feature telemetry windows
        onto a five-class taxonomy with drift detection;
  \item the \textbf{\jcac{} planner} (\emph{Joint Cache--Autoscaling
        Control}): receding-horizon model-predictive control over a
        move-blocked action lattice, minimizing
        $J = \alpha\,\text{cost} + \beta\,\text{SLO violation} +
        \gamma\,(1-\text{Jain})$, deployed as a Python service that imports
        the simulator's controller so offline and online planning are the
        same code;
  \item a \textbf{Kubernetes operator} (Go) with four CRDs (Tenant,
        WorkloadProfile, Policy, Budget) that actuates plans under
        guardrails: replica bounds, one cache level per step, per-tenant
        budget caps, and last-good-plan fallback.
\end{itemize}

\section{Contributions}

\begin{enumerate}
  \item \textbf{Joint cross-layer control.} The first surveyed controller to
        co-optimize replicas, semantic-cache budget, and model tier against
        one multi-tenant objective (Chapter~\ref{ch:design}); ablating
        jointness worsens the objective's cost term --- by $+2884\%$ against
        the published layered comparator, and by a margin $86\%$ smaller
        once that comparator's absorbing-tier defect is repaired
        (Section~\ref{sec:eval-v1}); the smaller number is the claim.
  \item \textbf{A measured composition result, and its adjudication.} On a
        1{,}800-run blocked factorial with tuned baselines, \polyforge{} wins
        the composite objective against every baseline ($|\dz|$ 0.66--1.44)
        and is the only Pareto-undominated system --- non-dominated in 48 of
        60 cells under \emph{any} weighting; the composite-$J$ win holds in
        every later campaign that tested it, including against an
        offline-trained learned joint controller over the same action space.
        The \emph{cost} component was then re-scored against comparators
        corrected for three confounds found after the campaigns closed
        (Section~\ref{sec:eval-adjudication}): the replayed-BurstGPT $-70\%$
        and Azure $-42\%$ headlines are \textbf{withdrawn} as like-for-like
        comparisons, and what survives is a feasibility result --- the
        per-tenant budget is satisfiable only by a controller holding the
        tier knob, and no replica-only reactive controller can meet it under
        AI load by any decision available to it.
  \item \textbf{An honest SLO result.} Two pre-registered attempts to beat
        tuned reactive scalers on raw violation failed and are published as
        failures; what holds is violation parity at $-43$\ldots$-48\%$ cost,
        a confirmed forecasting mechanism (seasonal $<$ trend on violation
        where real periodicity exists, $p=6\times10^{-4}$), and a
        declared-exploratory win in the spike regime where reaction is
        structurally impossible (Section~\ref{sec:eval-v3}). On the
        unbounded severity metric the synthetic and Azure advantage is real
        and the BurstGPT one reverses; both are stated together, always.
  \item \textbf{Real-data cache results with a quality denominator.} On real
        LMSYS-Chat-1M prompts: 29.6\% adaptive hit rate at cosine 0.85
        ($+93\%$ over a 10\%-of-working-set fixed cache), dollar-weighted
        savings, an empirical hit-rate--capacity curve, and --- unique among
        surveyed cache papers --- measured hit \emph{precision} with its
        stochasticity ceiling (Section~\ref{sec:eval-cache}).
  \item \textbf{Fairness at the control-plane altitude.} Beats a FIRM-style
        baseline on Jain's index ($\dz{=}1.0$) and a VTC-style replica-level
        transplant \emph{on Jain itself} at $-35\%$ cost; the controller's
        own fairness term is honestly nulled under injected interference
        (Section~\ref{sec:eval-fairness}).
  \item \textbf{A security frontier for shared semantic caches.} A timing
        side channel leaking prompt membership at AUC 0.88, and a defense
        comparison in which per-tenant partitioning restores
        attacker-to-chance at 24\% latency cost, dominating padding and TTL
        jitter (Section~\ref{sec:eval-security}).
  \item \textbf{Construction-level theory.} Block-coordinate optimality of
        the planner's two-sweep solver and boundedness/contraction of its
        self-calibration scale, stated at exactly the altitude the code
        delivers (Section~\ref{sec:design-theory}); a single-tenant cost
        separation between reactive and predictive control computed from
        the plant constants alone, shown to be exactly the information
        asymmetry of observational aliasing; and its multi-tenant extension
        computed exactly and \emph{falsified} against the measured arms
        (Section~\ref{sec:eval-summary}).
  \item \textbf{Live evidence at three depths.} A replica-only ordinal check
        on a kind cluster (disagreement, published), a valid 24-hour soak
        (25.8M requests, nothing dropped, two of four hypotheses failed and
        published), a real trace driving the real cluster, and a three-knob
        live plane on a rented GPU host on which the published controller
        \emph{lost} to its own ablation, the corrected controller then beat
        every arm at iso-fairness, and the obvious damper for its
        oscillation failed (Section~\ref{sec:eval-live}).
  \item \textbf{A pre-registered systems methodology.} Forty-eight
        protocols pushed before their runs, tuned baselines, paired effect
        sizes with bootstrap CIs, stopping rules, substrate separation, an
        adjudication practice in which confounds found after a campaign
        closes are answered by a new pre-registered record and never by
        editing the old one, and a published honest-nulls ledger
        (Chapter~\ref{ch:methodology}) --- a methodological contribution
        independent of any performance number.
\end{enumerate}

\section{Scope and honesty of claims}
\label{sec:intro-scope}

Every headline number in this thesis is \emph{decision quality under a stated
system model}: a discrete-event simulator whose congestion form was
calibrated against a real inference server (the fit direction errs
\emph{against} \polyforge{}'s lean postures) and whose model-tier constants
were checked against a real GPU latency bench. Real data enters three ways:
replayed BurstGPT demand (10.63M requests) and Azure LLM inference 2024
demand (44.1M requests) for the cost and forecasting claims, and real
LMSYS-Chat-1M prompts for the cache claims. Live-cluster validation never
compares absolutes across substrates. The replica-only ordinal check
recorded a \emph{disagreement} (live parity between the arms), published as
measured; the three-knob live plane that followed on a rented GPU host is
the one place where every knob is live and every arm is scored on the same
cells, and its verdicts --- the published controller's loss, the corrected
controller's win, the damper's failure --- are reported in
Section~\ref{sec:eval-live} exactly as their frozen protocols scored them.
What is \emph{not} claimed: any comparative cost percentage against the
originally tuned baselines (withdrawn, Section~\ref{sec:eval-adjudication}),
absolute dollar or latency numbers transferring to any production fleet,
token-level scheduling fairness, or tail behavior beyond p95 on the
simulated substrate.

\section{Document structure}

Chapter~\ref{ch:background} positions the work against autoscaling, caching,
fairness, and the 2026 serving stack. Chapter~\ref{ch:design} presents the
architecture, the \jcac{} formulation, and the two propositions.
Chapter~\ref{ch:implementation} covers the implementation and its
engineering gates. Chapter~\ref{ch:methodology} defines the pre-registered
methodology. Chapter~\ref{ch:evaluation} reports all campaigns ---
wins, nulls, and failures. Chapter~\ref{ch:discussion} confronts threats to
validity. Chapter~\ref{ch:conclusion} concludes.
Appendix~\ref{app:traceability} maps every claim to its committed evidence.
